\section{Introdução: Além do Horizonte Biomecânico}

Os gêmeos digitais na odontologia encontram-se no limiar de uma transição epocal: a evolução de protótipos acadêmicos para a adoção clínica sistêmica. A trajetória percorrida ao longo desta obra revelou uma ontologia em rápida maturação, desde os seus fundamentos computacionais até a consolidação de ecossistemas orquestrados como o OpenCode (\autoref{ch:opencode}). A Parte IV, em particular, desnudou os imperativos bioéticos (\autoref{ch:etica}) e as barreiras infraestruturais (\autoref{ch:desafios}) que modulam a absorção tecnológica no Brasil e no mundo.

Este capítulo final não se propõe a ser um epílogo especulativo, mas sim um \textit{roadmap} prospectivo. O objetivo é projetar cenários ancorados nas trajetórias das patentes atuais e na convergência de tecnologias emergentes — como a fusão multimodal e a modelagem generativa — delineando uma agenda rigorosa de pesquisa e desenvolvimento (P\&D). Ao transpor o horizonte puramente biomecânico rumo à saúde preditiva de base genômica e populacional, o texto orienta a comunidade científica, o ecossistema HealthTech e os \textit{policymakers} do Sistema Único de Saúde para a próxima década de inovações odontológicas in silico.
